% ==============================================================================
% Chapter 7: Unsupervised Learning & Cluster Analysis
% ==============================================================================
\chapter{Unsupervised Learning \& Cluster Analysis}
\label{ch:unsupervised_and_clustering}

\section{Foundations of Unsupervised Learning}

\subsection{Conceptual Framework and Objectives}
\textbf{Cluster Analysis} comprises a foundational family of unsupervised machine learning techniques designed to group unlabeled observations such that objects assigned to the same group (termed a \textbf{cluster}) exhibit significantly greater mutual similarity than objects assigned to alternate clusters.

Unlike supervised learning paradigms where ground-truth target labels guide functional mapping, unsupervised learning operates on unlabeled input datasets:
\begin{itemize}
    \item \textbf{Raw Data Input:} Observations $\mathbf{x}_i \in \mathbb{R}^p$ possess no associated response variables or class labels ($y_i$ is entirely absent).
    \item \textbf{Discovered Clusters:} The learning algorithm autonomously extracts and reveals the intrinsic latent geometric and statistical structure of the feature space.
    \item \textbf{Optimization Objective:} Formally, clustering seeks to simultaneously \textbf{maximize cohesion} (ensuring observations within the identical cluster reside close together) and \textbf{maximize separation} (ensuring distinct clusters remain far apart in metric space). Mathematically, this corresponds to:
    \begin{equation}
        \min \sum_{k=1}^K \sum_{\mathbf{x} \in C_k} d(\mathbf{x}, \boldsymbol{\mu}_k) \quad \text{and} \quad \max_{j \neq k} d(C_j, C_k)
    \end{equation}
    minimizing within-cluster dispersion while maximizing between-cluster distances.
\end{itemize}

\subsection{Supervised versus Unsupervised Learning}
The fundamental conceptual divergence between supervised classification and unsupervised clustering is summarized below:

\begin{table}[htbp]
\centering
\caption{Supervised Classification versus Unsupervised Clustering Paradigm}
\label{tab:supervised_vs_unsupervised}
\vspace{0.2em}
\renewcommand{\arraystretch}{1.3}
\begin{tabularx}{\textwidth}{>{\bfseries}lXX}
\toprule
Dimension & Supervised Learning (Classification) & Unsupervised Learning (Clustering) \\
\midrule
Target Labels & Training instances are paired with explicit ground-truth response labels ($y$). & No target labels or response variables exist in the dataset. \\
Primary Goal & Learn an inductive functional mapping $f(\mathbf{x}) \to y$ to predict unseen labels. & Permit the underlying data geometry to partition observations into natural groupings. \\
Core Question & \textit{``What is the predicted label of this new observation?''} & \textit{``Which observations naturally belong together?''} \\
Feedback Signal & Direct error feedback via supervised loss functions ($y - \hat{y}$). & Indirect structural feedback via cluster compactness and separation. \\
\bottomrule
\end{tabularx}
\end{table}

\subsection{The Standard Clustering Pipeline}
Executing cluster analysis on complex datasets follows an end-to-end, multi-stage operational workflow:
\begin{enumerate}
    \item \textbf{Data Preparation:} Clean raw features, manage missing values, and rescale heterogeneous dimensions (e.g., via $Z$-score standardization or Min-Max normalization).
    \item \textbf{Proximity Computation:} Select an appropriate metric distance or similarity measure aligned with feature types.
    \item \textbf{Algorithm Selection:} Choose a clustering paradigm suitable for dataset geometry (e.g., Hierarchical, $K$-Means, DBSCAN, or Gaussian Mixture Models).
    \item \textbf{Cluster Formation:} Execute the clustering algorithm to partition the feature space or build dendrogram hierarchies.
    \item \textbf{Cluster Cardinality Determination:} Evaluate the optimal number of clusters ($k$) via objective geometric heuristics.
    \item \textbf{Interpretation \& Profiling:} Inspect, label, profile, and economically interpret the characteristic centroids of each discovered group.
\end{enumerate}

\subsection{The Data Matrix and Variable Typology}
The input dataset is formalized as an $n \times p$ data matrix $\mathbf{X}$, where $n$ denotes the total number of objects (rows) and $p$ denotes the dimensionality of measured features (columns):
\begin{equation}
\mathbf{X} = \begin{bmatrix}
x_{11} & x_{12} & \cdots & x_{1p} \\
x_{21} & x_{22} & \cdots & x_{2p} \\
\vdots & \vdots & \ddots & \vdots \\
x_{n1} & x_{n2} & \cdots & x_{np}
\end{bmatrix} \in \mathbb{R}^{n \times p}
\end{equation}

Selecting a mathematically rigorous proximity measure requires classifying features according to their scale of measurement:

\begin{table}[!htbp]
\centering
\caption{Taxonomy of Data Types, Empirical Properties, and Proximity Measures}
\label{tab:cluster_variable_types}
\vspace{0.2em}
\renewcommand{\arraystretch}{1.3}
\begin{tabularx}{\textwidth}{>{\bfseries}l>{\raggedright\arraybackslash}p{2.8cm}XX}
\toprule
Type & Mathematical Meaning & Quantitative Finance Example & Typical Proximity Metric \\
\midrule
Nominal & Categorical states with no inherent natural ordering. & Industrial sector: Banking, IT, Energy, Healthcare. & Matching coefficient (Simple coincidence: same vs. different). \\
Binary & Two-state indicator variable taking values in $\{0, 1\}$. & Loan default event ($\text{Yes} = 1, \text{No} = 0$). & Simple Matching Coefficient (SMC), Jaccard similarity. \\
Ordinal & Categorical states possessing a distinct hierarchical order. & Credit rating: $\text{AAA} > \text{AA} > \text{A} > \text{BBB} > \dots$ & Normalized ranks mapped to the unit interval $[0, 1]$. \\
Interval & Equal incremental steps with an arbitrary (non-true) zero point. & Temperature in $^\circ\text{C}$, calendar dates. & Euclidean distance ($L_2$), Manhattan distance ($L_1$). \\
Ratio & Equal incremental steps anchored by an absolute, non-arbitrary zero. & Asset price, trading volume, portfolio return, leverage ratio. & Euclidean distance ($L_2$), Manhattan distance ($L_1$). \\
\bottomrule
\end{tabularx}
\end{table}

% ------------------------------------------------------------------------------
\section{Proximity, Distance Metrics, and Metric Axioms}

\subsection{Similarity, Dissimilarity, and Conversion}
Proximity is the overarching mathematical term denoting either mutual closeness or mutual divergence between two feature vectors:
\begin{itemize}
    \item \textbf{Similarity ($s$):} A quantitative measure of how alike two objects are. Higher values indicate greater proximity, typically constrained within the unit interval:
    \begin{equation}
        s(\mathbf{p}, \mathbf{q}) \in [0, 1]
    \end{equation}
    \item \textbf{Dissimilarity ($d$):} A quantitative measure of how different two objects are. Lower values indicate greater proximity, with a theoretical minimum of zero:
    \begin{equation}
        d(\mathbf{p}, \mathbf{q}) \ge 0, \quad d(\mathbf{p}, \mathbf{p}) = 0
    \end{equation}
    \item \textbf{Standard Conversion:} When dissimilarity is bounded within the unit interval $d \in [0, 1]$, an equivalent similarity measure is defined via the complementary transformation:
    \begin{equation}
        s = 1 - d
    \end{equation}
\end{itemize}

\subsection{Formal Metric Axioms for Distance Functions}
A distance function $d(\mathbf{p}, \mathbf{q})$ on a metric space must satisfy four fundamental mathematical axioms for all points $\mathbf{p}, \mathbf{q}, \mathbf{r}$:
\begin{enumerate}
    \item \textbf{Non-Negativity:} The distance between any two distinct observations is strictly non-negative:
    \begin{equation}
        d(\mathbf{p}, \mathbf{q}) \ge 0
    \end{equation}
    \item \textbf{Identity of Indiscernibles:} The distance between two points is zero if and only if the points are identical:
    \begin{equation}
        d(\mathbf{p}, \mathbf{q}) = 0 \iff \mathbf{p} = \mathbf{q}
    \end{equation}
    \item \textbf{Symmetry:} The distance from $\mathbf{p}$ to $\mathbf{q}$ equals the distance from $\mathbf{q}$ to $\mathbf{p}$:
    \begin{equation}
        d(\mathbf{p}, \mathbf{q}) = d(\mathbf{q}, \mathbf{p})
    \end{equation}
    \item \textbf{Triangle Inequality:} The direct distance between two points cannot exceed the sum of indirect distances through any intermediate point:
    \begin{equation}
        d(\mathbf{p}, \mathbf{r}) \le d(\mathbf{p}, \mathbf{q}) + d(\mathbf{q}, \mathbf{r})
    \end{equation}
\end{enumerate}

\begin{figure}[htbp]
\centering
\begin{tikzpicture}[scale=1.1]
    % Coordinates of metric triangle
    \coordinate (P) at (0, 0);
    \coordinate (Q) at (2.2, 2.5);
    \coordinate (R) at (5.0, 0);

    % Draw metric edges
    \draw[navyblue, very thick] (P) -- node[left=4pt, font=\small\sffamily, text=navyblue] {$d(\mathbf{p}, \mathbf{q})$} (Q);
    \draw[navyblue, very thick] (Q) -- node[right=4pt, font=\small\sffamily, text=navyblue] {$d(\mathbf{q}, \mathbf{r})$} (R);
    \draw[darkteal, line width=1.4pt] (P) -- node[below=4pt, font=\small\sffamily, text=darkteal] {$d(\mathbf{p}, \mathbf{r}) \le d(\mathbf{p}, \mathbf{q}) + d(\mathbf{q}, \mathbf{r})$} (R);

    % Draw vertices
    \filldraw[crimson] (P) circle (2.5pt) node[below left=2pt, font=\bfseries] {$\mathbf{p}$};
    \filldraw[crimson] (Q) circle (2.5pt) node[above=3pt, font=\bfseries] {$\mathbf{q}$};
    \filldraw[crimson] (R) circle (2.5pt) node[below right=2pt, font=\bfseries] {$\mathbf{r}$};
\end{tikzpicture}
\caption{Geometric schematic of the metric triangle inequality $d(\mathbf{p}, \mathbf{r}) \le d(\mathbf{p}, \mathbf{q}) + d(\mathbf{q}, \mathbf{r})$.}
\label{fig:triangle_inequality}
\end{figure}

\subsection{Standard Continuous Distance Metrics}
For continuous real-valued vectors $\mathbf{p}, \mathbf{q} \in \mathbb{R}^p$, metric distance is evaluated using standard norms:
\begin{enumerate}
    \item \textbf{Euclidean Distance ($L_2$ Norm):} Measures straight-line geometric distance:
    \begin{equation}
        d(\mathbf{p}, \mathbf{q}) = \|\mathbf{p} - \mathbf{q}\|_2 = \sqrt{\sum_{k=1}^p (p_k - q_k)^2}
    \end{equation}
    \item \textbf{Manhattan Distance ($L_1$ Norm / City-Block Metric):} Measures rectilinear grid traversal distance:
    \begin{equation}
        d(\mathbf{p}, \mathbf{q}) = \|\mathbf{p} - \mathbf{q}\|_1 = \sum_{k=1}^p |p_k - q_k|
    \end{equation}
    \item \textbf{Minkowski Distance ($L_r$ Norm):} A generalized metric parameterized by order $r \ge 1$:
    \begin{equation}
        d_r(\mathbf{p}, \mathbf{q}) = \left( \sum_{k=1}^p |p_k - q_k|^r \right)^{1/r}
    \end{equation}
    Notice that $r = 1$ recovers the Manhattan metric, and $r = 2$ recovers the Euclidean metric.
    \item \textbf{Mahalanobis Distance:} Accounts for scale variance and inter-feature covariance:
    \begin{equation}
        d_M(\mathbf{p}, \mathbf{q}) = \sqrt{(\mathbf{p} - \mathbf{q})^T \boldsymbol{\Sigma}^{-1} (\mathbf{p} - \mathbf{q})}
    \end{equation}
    where $\boldsymbol{\Sigma} \in \mathbb{R}^{p \times p}$ denotes the feature covariance matrix. When features are uncorrelated with unit variance ($\boldsymbol{\Sigma} = \mathbf{I}$), Mahalanobis distance simplifies to standard Euclidean distance.
\end{enumerate}

\subsection{Proximity for Binary and Mixed Data}
When evaluating binary indicator vectors, co-occurrence contingency frequencies are categorized into four mutual states:
\begin{itemize}
    \item $a = $ count of features where both object $i$ and object $j$ possess value $1$ ($\{i=1, j=1\}$).
    \item $b = $ count of features where object $i$ has $1$ and object $j$ has $0$ ($\{i=1, j=0\}$).
    \item $c = $ count of features where object $i$ has $0$ and object $j$ has $1$ ($\{i=0, j=1\}$).
    \item $d = $ count of features where both object $i$ and object $j$ possess value $0$ ($\{i=0, j=0\}$).
\end{itemize}

\begin{enumerate}
    \item \textbf{Simple Matching Coefficient (SMC):} Weights shared presence ($1$-$1$) and shared absence ($0$-$0$) symmetrically:
    \begin{equation}
        \text{SMC} = \frac{a + d}{a + b + c + d}
    \end{equation}
    \item \textbf{Jaccard Similarity Coefficient:} Ignores shared negative matches ($d$), focusing strictly on mutual active attributes:
    \begin{equation}
        J = \frac{a}{a + b + c}
    \end{equation}
    This formulation is essential when $0$ reflects the passive absence of an attribute (e.g., non-transaction or unpurchased asset).
    \item \textbf{Gower Distance:} For heterogeneous datasets containing a combination of continuous, ordinal, and nominal variables, Gower's metric computes a normalized per-variable distance $d_k(i, j) \in [0, 1]$ and computes their weighted average across all valid attributes:
    \begin{equation}
        d_{\text{Gower}}(i, j) = \frac{\sum_{k=1}^p w_k d_k(i, j)}{\sum_{k=1}^p w_k}
    \end{equation}
\end{enumerate}

% ------------------------------------------------------------------------------
\section{Taxonomy of Clustering Methodologies}

Clustering algorithms are systematically categorized into four foundational methodological families:

\begin{figure}[!htbp]
\centering
\begin{tikzpicture}[
    root/.style={
        rectangle, draw=navyblue, fill=navyblue!10, rounded corners=3pt,
        line width=1.2pt, minimum width=3.4cm, minimum height=0.9cm, align=center,
        font=\bfseries\sffamily\small, text=navyblue
    },
    branch/.style={
        rectangle, draw=navyblue, fill=lightgraybg, rounded corners=2pt,
        thick, minimum width=2.7cm, minimum height=1.0cm, align=center,
        font=\bfseries\sffamily\footnotesize, text=navyblue
    },
    leaf/.style={
        rectangle, draw=darkteal, fill=darkteal!8, rounded corners=2pt,
        thick, minimum width=2.4cm, minimum height=0.9cm, align=center,
        font=\bfseries\sffamily\scriptsize, text=darkteal
    },
    edge/.style={draw=navyblue, thick, -{Stealth[length=2.2mm]}},
    fork/.style={draw=darkteal, thick, -{Stealth[length=2.0mm]}}
]
    % Root
    \node[root] (root) at (0, 2.2) {Clustering Methods};

    % 4 Branches
    \node[branch] (b1) at (-5.1, 0.5) {Partition-Based\\(e.g., $K$-Means)};
    \node[branch] (b2) at (-1.7, 0.5) {Hierarchical\\Clustering};
    \node[branch] (b3) at (1.7, 0.5) {Density-Based\\(e.g., DBSCAN)};
    \node[branch] (b4) at (5.1, 0.5) {Model / Graph\\Based (e.g., GMM)};

    % Connect root to 4 branches with clean tree fork
    \coordinate (rfork) at (0, 1.35);
    \draw[navyblue, thick] (root.south) -- (rfork);
    \draw[edge] (rfork) -| (b1.north);
    \draw[edge] (rfork) -| (b2.north);
    \draw[edge] (rfork) -| (b3.north);
    \draw[edge] (rfork) -| (b4.north);

    % Leaves under Hierarchical with clear 0.6cm clearance and drop-stems
    \node[leaf] (l1) at (-3.2, -1.3) {Agglomerative\\(Bottom-Up)};
    \node[leaf] (l2) at (-0.2, -1.3) {Divisive\\(Top-Down)};

    \coordinate (hfork) at (-1.7, -0.4);
    \draw[darkteal, thick] (b2.south) -- (hfork);
    \draw[fork] (hfork) -| (l1.north);
    \draw[fork] (hfork) -| (l2.north);
\end{tikzpicture}
\caption{Taxonomy of major unsupervised clustering paradigms.}
\label{fig:clustering_taxonomy}
\end{figure}

% ------------------------------------------------------------------------------
\section{Hierarchical Clustering}

\subsection{Dendrogram Architecture and Horizontal Cuts}
\textbf{Hierarchical Clustering} constructs a continuous nested tree of partitions visualized as a \textbf{dendrogram}:
\begin{itemize}
    \item \textbf{Leaves:} Individual ground observations at the baseline ($h = 0$).
    \item \textbf{Merges / Splits:} Merges occurring at lower dissimilarity heights indicate strong mutual affinity; merges occurring at higher heights indicate substantial dissimilarity.
    \item \textbf{Horizontal Threshold Cutting:} A horizontal slice drawn across the dendrogram at height $h^*$ defines the final number of isolated clusters. Lowering the cut threshold increases the number of fine-grained clusters; raising the cut threshold yields fewer, coarser clusters.
\end{itemize}

\begin{figure}[htbp]
\centering
\begin{tikzpicture}[x=1.3cm, y=1.1cm]
    % Vertical dissimilarity axis
    \draw[->, >=stealth, thick, darkgray] (0, 0) -- (0, 4.8) node[above, font=\small\sffamily] {Dissimilarity ($h$)};

    % Height ticks
    \foreach \y/\label in {1/h_1, 2/h_2, 3/h_3, 4/h_4} {
        \draw (-0.08, \y) -- (0.08, \y) node[left=4pt, font=\footnotesize] {$\label$};
    }

    % Dendrogram branches (observations 1 to 5)
    % Merge (1, 2) at h1
    \draw[navyblue, thick] (1, 0) -- (1, 1) -- (2, 1) -- (2, 0);
    % Merge (3, 4) at h1
    \draw[navyblue, thick] (3.5, 0) -- (3.5, 1) -- (4.5, 1) -- (4.5, 0);
    % Merge ((3,4), 5) at h2
    \draw[navyblue, thick] (4, 1) -- (4, 2) -- (5.5, 2) -- (5.5, 0);
    % Merge ((1,2), ((3,4),5)) at h3
    \draw[navyblue, thick] (1.5, 1) -- (1.5, 3) -- (4.75, 3) -- (4.75, 2);
    % Top root at h4
    \draw[navyblue, thick] (3.125, 3) -- (3.125, 4.2);

    % Horizontal cut lines
    \draw[dashed, crimson, thick] (0.2, 2.5) -- (6.0, 2.5) node[right, font=\footnotesize\bfseries, text=crimson] {Cut at $h^* \implies 2$ Clusters};
    \draw[dashed, darkteal, thick] (0.2, 1.5) -- (6.0, 1.5) node[right, font=\footnotesize\bfseries, text=darkteal] {Cut at $h^{**} \implies 3$ Clusters};

    % Leaf nodes
    \node[below=2pt, font=\scriptsize\bfseries] at (1, 0) {$\mathbf{x}_1$};
    \node[below=2pt, font=\scriptsize\bfseries] at (2, 0) {$\mathbf{x}_2$};
    \node[below=2pt, font=\scriptsize\bfseries] at (3.5, 0) {$\mathbf{x}_3$};
    \node[below=2pt, font=\scriptsize\bfseries] at (4.5, 0) {$\mathbf{x}_4$};
    \node[below=2pt, font=\scriptsize\bfseries] at (5.5, 0) {$\mathbf{x}_5$};
\end{tikzpicture}
\caption{Hierarchical clustering dendrogram illustrating horizontal cut thresholds determining cluster cardinality.}
\label{fig:dendrogram_structure}
\end{figure}

\subsection{Agglomerative versus Divisive Paradigms}
Hierarchical tree construction operates under two alternative procedural philosophies:
\begin{enumerate}
    \item \textbf{Agglomerative (Bottom-Up):} Initializes each observation in its own singleton cluster ($n$ initial clusters). The algorithm iteratively locates and merges the two closest clusters until a single grand cluster contains all points. Agglomerative approaches are computationally efficient ($O(n^2)$ space, $O(n^3)$ worst-case time) and represent the standard practical choice.
    \item \textbf{Divisive (Top-Down):} Initializes all $n$ observations in one global cluster. The algorithm iteratively partitions clusters into sub-clusters until every observation is isolated in a singleton cluster. Because evaluating all candidate bipartitions exhibits exponential complexity ($2^{n-1} - 1$), unconstrained divisive clustering is computationally prohibitive and rarely used without greedy splitting heuristics.
\end{enumerate}

\subsection{The Agglomerative Induction Algorithm}
The procedural steps for agglomerative clustering are formalized as follows:
\begin{enumerate}
    \item Compute the pairwise distance matrix $\mathbf{D} \in \mathbb{R}^{n \times n}$ between all $n$ initial observations.
    \item Assign each object $\mathbf{x}_i$ to its own singleton cluster $C_i = \{\mathbf{x}_i\}$.
    \item Identify the two clusters $A$ and $B$ that minimize inter-cluster dissimilarity $d(A, B)$.
    \item Merge clusters $A$ and $B$ into a composite cluster $A \cup B$, and record the dissimilarity value as the branch merge height.
    \item Update the distance matrix $\mathbf{D}$ by evaluating distances between the newly formed cluster $A \cup B$ and all remaining clusters using a pre-specified \textbf{linkage rule}.
    \item Repeat Steps 3--5 until exactly one comprehensive cluster remains (requiring exactly $n - 1$ sequential merges).
\end{enumerate}

\subsection{Inter-Cluster Linkage Criteria}
The metric defining the distance between two multi-point clusters $A$ and $B$ governs dendrogram topology:

\begin{table}[!htbp]
\centering
\caption{Mathematical Definitions and Geometric Properties of Linkage Rules}
\label{tab:linkage_rules}
\vspace{0.2em}
\renewcommand{\arraystretch}{1.3}
\begin{tabularx}{\textwidth}{>{\bfseries\raggedright\arraybackslash}p{3.4cm}>{\raggedright\arraybackslash}p{4.0cm}X}
\toprule
Linkage Rule & Mathematical Definition & Geometric Characteristics \& Failure Modes \\
\midrule
Single Linkage (Min) & $d(A, B) = \min_{\mathbf{x} \in A, \mathbf{y} \in B} d(\mathbf{x}, \mathbf{y})$ & Optimistic measure evaluating the closest point pair. Highly prone to the \textbf{chaining effect}, forming extended, straggly clusters connected by single noise points. \\
Complete Linkage (Max) & $d(A, B) = \max_{\mathbf{x} \in A, \mathbf{y} \in B} d(\mathbf{x}, \mathbf{y})$ & Pessimistic measure evaluating the most distant point pair. Enforces compact, spherical clusters of roughly equal diameter; sensitive to extreme outliers. \\
Average Linkage (Mean) & $d(A, B) = \frac{1}{|A||B|} \sum_{\mathbf{x} \in A} \sum_{\mathbf{y} \in B} d(\mathbf{x}, \mathbf{y})$ & Evaluates the average pairwise distance between all member points. Robust compromise between single and complete linkage. \\
Centroid Linkage & $d(A, B) = \|\boldsymbol{\mu}_A - \boldsymbol{\mu}_B\|_2$ & Measures Euclidean distance between cluster centroids $\boldsymbol{\mu}_A, \boldsymbol{\mu}_B$. May produce non-monotonic dendrograms (tree inversions). \\
Ward's Minimum Variance & $\Delta \text{ESS} = \frac{|A||B|}{|A| + |B|} \|\boldsymbol{\mu}_A - \boldsymbol{\mu}_B\|^2$ & Merges the cluster pair minimizing the increase in total within-cluster sum of squared errors ($W(C) = \sum_{\mathbf{x} \in C} \|\mathbf{x} - \boldsymbol{\mu}_C\|^2$). Tends to find highly cohesive, spherical clusters. \\
\bottomrule
\end{tabularx}
\end{table}

\clearpage
\subsection{Divisive Splitting Strategies: Monothetic versus Polythetic}
Divisive hierarchical methods decompose clusters according to either single or multiple attributes:

\begin{table}[!htbp]
\centering
\caption{Monothetic versus Polythetic Divisive Clustering Characteristics}
\label{tab:monothetic_vs_polythetic}
\vspace{0.2em}
\renewcommand{\arraystretch}{1.3}
\begin{tabularx}{\textwidth}{>{\bfseries}lXX}
\toprule
Dimension & Monothetic Divisive Clustering & Polythetic Divisive Clustering \\
\midrule
Variables per Split & Evaluates exactly \textbf{one variable} at a time to partition clusters. & Evaluates \textbf{multiple variables simultaneously} to establish the split. \\
Splitting Basis & Based on a single threshold or categorical attribute match. & Based on multidimensional metric distance or dissimilarity. \\
Complexity & Computationally simpler, interpretable decision-tree-like splits. & More comprehensive, accounting for multi-attribute interactions. \\
Primary Algorithms & Decision tree splitting heuristics, binary attribute splits. & DIANA (Divisive Analysis), Farthest-Seed bipartitioning. \\
\bottomrule
\end{tabularx}
\end{table}

\subsection{Strengths, Limitations, and Empirical Heuristics}
Hierarchical clustering displays distinct operational tradeoffs:
\begin{itemize}
    \item \textbf{Strengths:}
    \begin{enumerate}
        \item Does not require pre-specifying the number of clusters $k$ in advance.
        \item The resulting dendrogram exposes multiscale nested structures across all levels of granularity.
        \item Deterministic algorithm: identical inputs and distance matrices guarantee identical tree topologies.
        \item Compatible with any arbitrary dissimilarity matrix (including non-Euclidean proximity).
    \end{enumerate}
    \item \textbf{Limitations:}
    \begin{enumerate}
        \item Memory-intensive ($O(n^2)$ distance storage) and computationally demanding ($O(n^3)$ unoptimized time complexity).
        \item Merges and splits are strictly irrevocable: greedy local errors made in initial steps cannot be corrected.
        \item Highly sensitive to user-chosen linkage criteria.
        \item Dendrogram visualization degrades rapidly on large datasets ($n > 10{,}000$).
    \end{enumerate}
    \item \textbf{Thumb Rule:} Apply hierarchical clustering primarily to small-to-medium datasets (up to a few thousand objects) where exploring nested taxonomic hierarchies is paramount (e.g., biological taxonomy, credit rating transition classes, product brand segmentation).
\end{itemize}

% ------------------------------------------------------------------------------
\section{Non-Hierarchical Partitioning: The \texorpdfstring{$K$-Means}{K-Means} Algorithm}

\subsection{Foundational Principles and Procedure}
\textbf{Non-hierarchical partitioning} divides an unlabeled dataset of $n$ observations into a pre-specified number $K$ of non-overlapping, mutually exclusive clusters $C_1, C_2, \dots, C_K$.

\textbf{$K$-Means} is the preeminent partitioning algorithm, operating through an iterative Lloyd-Forgy optimization scheme minimizing the total Within-Cluster Sum of Squares (WCSS):
\begin{equation}
    \text{WCSS} = \sum_{j=1}^K \sum_{\mathbf{x}_i \in C_j} \|\mathbf{x}_i - \boldsymbol{\mu}_j\|^2
\end{equation}

The procedural execution pipeline proceeds through four sequential stages:
\begin{enumerate}
    \item \textbf{Initialization:} Select cluster cardinality $K$ in advance (e.g., $K = 3$). Initialize $K$ cluster centroids $\boldsymbol{\mu}_1, \boldsymbol{\mu}_2, \dots, \boldsymbol{\mu}_K \in \mathbb{R}^p$ (randomly or via $K$-Means++).
    \item \textbf{Assignment Step:} Assign each observation $\mathbf{x}_i$ to its nearest cluster centroid according to Euclidean distance:
    \begin{equation}
        c(i) = \arg\min_{j \in \{1, \dots, K\}} \|\mathbf{x}_i - \boldsymbol{\mu}_j\|_2 = \arg\min_{j \in \{1, \dots, K\}} \sqrt{\sum_{r=1}^p (x_{ir} - \mu_{jr})^2}
    \end{equation}
    \item \textbf{Update Step:} Recompute the centroid of each cluster $C_j$ as the component-wise empirical arithmetic mean of all observations currently assigned to it:
    \begin{equation}
        \boldsymbol{\mu}_j = \frac{1}{n_j} \sum_{\mathbf{x}_i \in C_j} \mathbf{x}_i \quad \implies \quad \mu_{jr} = \frac{1}{n_j} \sum_{i=1}^{n_j} x_{ir}
    \end{equation}
    where $n_j = |C_j|$ represents the number of observations belonging to cluster $C_j$.
    \item \textbf{Convergence \& Termination:} Repeat the assignment and update steps iteratively until:
    \begin{itemize}
        \item Centroid coordinates remain stable ($\|\boldsymbol{\mu}_j^{(t)} - \boldsymbol{\mu}_j^{(t-1)}\| < \epsilon$),
        \item Cluster assignment memberships cease to change, or
        \item A pre-defined maximum iteration count ($T_{\max}$) is reached.
    \end{itemize}
\end{enumerate}

\begin{figure}[htbp]
\centering
\begin{tikzpicture}[
    node distance=1.2cm and 1.0cm, auto,
    proc/.style={
        rectangle, draw=navyblue, fill=navyblue!5, rounded corners=3pt,
        line width=1.1pt, minimum width=2.8cm, minimum height=1.0cm, align=center,
        font=\sffamily\bfseries\small, text=navyblue
    },
    dec/.style={
        diamond, draw=crimson, fill=crimson!8, aspect=1.8,
        line width=1.1pt, inner sep=3pt, align=center,
        font=\sffamily\bfseries\footnotesize, text=crimson
    },
    stop/.style={
        rectangle, draw=darkteal, fill=darkteal!10, rounded corners=3pt,
        line width=1.2pt, minimum width=2.4cm, minimum height=1.0cm, align=center,
        font=\sffamily\bfseries\small, text=darkteal
    },
    arr/.style={->, >=stealth, very thick, navyblue}
]
    \node[proc] (init) {1. Initialize $K$\\Centroids};
    \node[proc, right=1.0cm of init] (assign) {2. Assign Points to\\Nearest Centroid};
    \node[proc, right=1.0cm of assign] (update) {3. Recalculate Means\\(Update Centroids)};
    \node[dec, below=1.4cm of update] (conv) {Converged?\\($\Delta \boldsymbol{\mu} < \epsilon$)};
    \node[stop, left=2.0cm of conv] (done) {Final Stable\\Clusters};

    \draw[arr] (init) -- (assign);
    \draw[arr] (assign) -- (update);
    \draw[arr] (update) -- (conv);
    \draw[arr, darkteal] (conv) -- node[above, font=\footnotesize\bfseries, text=darkteal] {Yes} (done);
    \draw[->, >=stealth, very thick, crimson] (conv.east) -- ++(0.9cm, 0) |- ([yshift=0.6cm]assign.north) node[pos=0.75, above, font=\footnotesize\bfseries, text=crimson] {No (Reassign Points)} -- (assign.north);
\end{tikzpicture}
\caption{Iterative assignment and centroid update cycle of the $K$-Means clustering algorithm.}
\label{fig:kmeans_pipeline}
\end{figure}

% ------------------------------------------------------------------------------
\section{Density-Based Clustering: DBSCAN, HDBSCAN, and OPTICS}

\subsection{Core Principles of Density Clustering}
Traditional centroid-based algorithms (such as $K$-Means) implicitly assume convex, spherically distributed clusters and exhibit severe degradation when confronted with arbitrary geometries, non-convex manifolds (e.g., concentric circles, interlocking spirals, S-shapes), or background noise.

\textbf{Density-Based Spatial Clustering of Applications with Noise (DBSCAN)} overcomes these limitations by defining clusters as continuous dense regions of points separated by sparse regions of empty space or noise. A cluster expands continuously along connected paths as long as the local neighborhood density exceeds a pre-defined threshold.

\subsection{Local Point Typology}
Density around a candidate observation $\mathbf{p}$ is evaluated within an $\epsilon$-neighborhood:
\begin{equation}
    N_\epsilon(\mathbf{p}) = \{\mathbf{q} \in D \mid \|\mathbf{p} - \mathbf{q}\| \le \epsilon\}
\end{equation}
where $\epsilon$ denotes the neighborhood radius and $\text{MinPts}$ denotes the minimum cardinality threshold. Data points are partitioned into three mutually exclusive categories:
\begin{enumerate}
    \item \textbf{Core Points:} Any point $\mathbf{p}$ possessing at least $\text{MinPts}$ neighbors within its $\epsilon$-radius:
    \begin{equation}
        |N_\epsilon(\mathbf{p})| \ge \text{MinPts}
    \end{equation}
    \item \textbf{Border Points:} A point that is not a core point ($|N_\epsilon(\mathbf{p})| < \text{MinPts}$), but falls directly within the $\epsilon$-neighborhood of at least one valid core point.
    \item \textbf{Noise Points (Outliers):} Any point that is neither a core point nor a border point. Noise points inhabit sparse background regions and are discarded from cluster assignments.
\end{enumerate}

\begin{figure}[!htbp]
\centering
\begin{tikzpicture}[scale=1.1,
    card/.style={
        rounded corners=3pt, align=center, font=\sffamily,
        minimum height=1.3cm, inner sep=4pt
    }
]
    % Top Symmetrical Cards with clear gaps and balanced 3 lines
    \node[card, text width=3.2cm, draw=navyblue, fill=navyblue!6, text=navyblue] (c_core) at (-2.3, 2.9) {
        \textbf{Core Point $\mathbf{p}$}\\[2pt]
        \scriptsize $|N_\epsilon(\mathbf{p})| = 6 \ge \text{MinPts} = 5$\\[1pt]
        \scriptsize (Dense Cluster Core)
    };
    \node[card, text width=3.2cm, draw=darkteal, fill=darkteal!6, text=darkteal] (c_border) at (1.4, 2.9) {
        \textbf{Border Point $\mathbf{b}$}\\[2pt]
        \scriptsize $|N_\epsilon(\mathbf{b})| = 3 < \text{MinPts}$\\[1pt]
        \scriptsize but $\mathbf{b} \in N_\epsilon(\mathbf{p})$
    };
    \node[card, text width=3.2cm, draw=crimson, fill=crimson!6, text=crimson] (c_noise) at (5.2, 2.9) {
        \textbf{Noise / Outlier $\mathbf{z}$}\\[2pt]
        \scriptsize $|N_\epsilon(\mathbf{z})| = 1 < \text{MinPts}$\\[1pt]
        \scriptsize (Unassigned Outlier)
    };

    % Coordinates for points (all centers on identical baseline y = 0.0)
    \coordinate (Core) at (-2.0, 0.0);
    \coordinate (Border) at (-0.7, 0.0);
    \coordinate (Noise) at (5.2, 0.0);

    % 1. Shaded backgrounds first (both transparent/light fills)
    \fill[navyblue!8] (Core) circle (1.6cm);
    \fill[darkteal!8] (Border) circle (1.6cm);
    \fill[crimson!8] (Noise) circle (1.6cm);

    % 2. Dashed boundaries drawn ON TOP so neither circle obscures the other
    \draw[dashed, navyblue, thick] (Core) circle (1.6cm);
    \draw[dashed, darkteal, thick] (Border) circle (1.6cm);
    \draw[dashed, crimson, thick] (Noise) circle (1.6cm);

    % Core point and its dense neighborhood (5 core cluster points + border point = 6)
    \filldraw[navyblue] (Core) circle (3.2pt) node[below=5pt, font=\footnotesize\bfseries\sffamily] {$\mathbf{p}$};
    \filldraw[navyblue] (-3.0, 0.7) circle (2.2pt);
    \filldraw[navyblue] (-2.8, -0.7) circle (2.2pt);
    \filldraw[navyblue] (-2.2, 1.0) circle (2.2pt);
    \filldraw[navyblue] (-2.1, -1.0) circle (2.2pt);

    % Border point (reachable from p, has only 3 points in its neighborhood: p, b, b1)
    \filldraw[darkteal] (Border) circle (3.2pt) node[below=5pt, font=\footnotesize\bfseries\sffamily] {$\mathbf{b}$};
    \filldraw[darkteal] (0.3, 0.5) circle (2.2pt);

    % Noise point (completely isolated outlier with no cluster points within epsilon)
    \filldraw[crimson] (Noise) circle (3.2pt) node[below=5pt, font=\footnotesize\bfseries\sffamily] {$\mathbf{z}$};

    % Radius annotations
    \draw[->, >=stealth, navyblue, thick] (Core) -- ++(140:1.6cm) node[midway, above left, font=\scriptsize\bfseries] {$\epsilon$};
    \draw[->, >=stealth, darkteal, thick] (Border) -- ++(35:1.6cm) node[midway, above right, font=\scriptsize\bfseries] {$\epsilon$};
    \draw[->, >=stealth, crimson, thick] (Noise) -- ++(45:1.6cm) node[midway, above right, font=\scriptsize\bfseries] {$\epsilon$};
\end{tikzpicture}
\caption{DBSCAN point classification into Core, Border, and Noise points ($\text{MinPts} = 5$).}
\label{fig:dbscan_points}
\end{figure}

\subsection{Advanced Density Algorithms: HDBSCAN and OPTICS}
Beyond standard DBSCAN, advanced density algorithms resolve variable-density challenges:
\begin{enumerate}
    \item \textbf{HDBSCAN (Hierarchical DBSCAN):} Converts DBSCAN into a hierarchical framework by extracting clusters across variable density thresholds, eliminating the sensitivity of tuning a single static $\epsilon$ parameter.
    \item \textbf{OPTICS (Ordering Points To Identify the Clustering Structure):} Produces a linear reachability distance plot that exposes clustering structure across multiple spatial scales simultaneously.
\end{enumerate}

\subsection{Practical Considerations in Unsupervised Pipelines}
When deploying unsupervised clustering pipelines to financial datasets, practitioners must address four critical engineering requirements:
\begin{enumerate}
    \item \textbf{Feature Rescaling:} Distance-based metrics are severely distorted by magnitude differences; standardize features prior to distance computation.
    \item \textbf{Metric Alignment:} Pair continuous, binary, and categorical attributes with compatible metrics (e.g., Gower or Jaccard).
    \item \textbf{Curse of Dimensionality:} As feature count $p$ expands, spatial sparsity causes all pairwise Euclidean distances to converge, degrading cluster quality.
    \item \textbf{Imputation Integrity:} Because unsupervised methods cannot compute distance vectors with missing values, robust statistical imputation must precede clustering.
\end{enumerate}
